\section{Análisis de Resultados}

Los resultados obtenidos indican que el caudal de diseño de 2 260 CFM satisface el requerimiento de 12 renovaciones de aire por hora para el volumen de 320 m$^3$. Este valor se encuentra dentro de los rangos recomendados por ASHRAE 170 para espacios de laboratorio y por el WHO Laboratory Biosafety Manual para instalaciones BSL-2, con margen de seguridad respecto al mínimo de 6 ACH.

La potencia teórica del ventilador, calculada en 0.44 kW para un escenario de presión total de 250 Pa, es consistente con el orden de magnitud esperado para un ventilador de impulsión de caudal medio en aplicaciones de ventilación de laboratorio. La selección de un motor de 1.0 HP proporciona un factor de seguridad adecuado para compensar pérdidas adicionales por ensuciamiento de filtros, variaciones de eficiencia y condiciones de operación no ideales.

La velocidad de impulsión de 8 m/s en la boca del ventilador se ubica dentro del rango recomendado de 6 a 12 m/s para ventiladores centrífugos y axiales de impulsión directa. Esta velocidad permite un ingreso de aire ordenado sin generar ruido excesivo ni corrientes locales molestas en el interior del recinto.

La velocidad de exfiltración de 3.0 m/s en las rejillas de salida representa un compromiso entre dos requerimientos contrapuestos. Por un lado, una velocidad suficientemente alta genera la resistencia necesaria para mantener la presión positiva deseada (+12.5 a +25 Pa). Por otro lado, una velocidad excesiva aumentaría el ruido y la velocidad de entrada del aire en las zonas cercanas a las rejillas. El valor seleccionado se encuentra dentro del rango aceptable de 2.5 a 4.0 m/s para aplicaciones de laboratorio.

La configuración de tres rejillas de exfiltración distribuye la salida del aire y favorece un flujo cruzado efectivo cuando se ubican en paredes opuestas o perpendiculares a la impulsión. Esto reduce la probabilidad de zonas de estancamiento y mejora la uniformidad de la presión en el recinto. El modelo CFD posterior deberá confirmar que el campo de presión relativa es positivo en todo el volumen del laboratorio y que no existen recirculaciones significativas.

Finalmente, el balance de masas entre la entrada y la salida del recinto se cierra con una discrepancia inferior al 0.1 \%, lo cual valida la coherencia dimensional y cinemática de los datos de contorno propuestos para el CFD.
